\documentclass[main.tex]{subfiles}
\begin{document}


\begin{register}{H}{I2C\_SCL\_LOW\_PERIOD\_REG}{0x0000}\label{regdesc:I2CSCLLOWPERIODREG}
\regfield{(reserved)}{18}{14}{000000000000000000}%
\regfield{I2C\_SCL\_LOW\_PERIOD}{14}{0}{00000000000000}%
\reglabel{Reset}\regnewline%
\begin{regdesc}\begin{reglist}
\item [I2C\_SCL\_LOW\_PERIOD]I2C 为 Master 时，该寄存器用于配置 SCL 时钟信号的低电平持续的 APB 时钟周期数。（读/写）
\end{reglist}\end{regdesc}
\end{register}


\begin{register}{H}{I2C\_CTR\_REG}{0x0004}\label{regdesc:I2CCTRREG}
\regfield{(reserved)}{24}{8}{000000000000000000000000}%
\regfield{I2C\_RX\_LSB\_FIRST}{1}{7}{0}%
\regfield{I2C\_TX\_LSB\_FIRST}{1}{6}{0}%
\regfield{I2C\_TRANS\_START}{1}{5}{0}%
\regfield{I2C\_MS\_MODE}{1}{4}{0}%
\regfield{(reserved)}{1}{3}{0}%
\regfield{I2C\_SAMPLE\_SCL\_LEVEL}{1}{2}{0}%
\regfield{I2C\_SCL\_FORCE\_OUT}{1}{1}{1}%
\regfield{I2C\_SDA\_FORCE\_OUT}{1}{0}{1}%
\reglabel{Reset}\regnewline%
\begin{regdesc}\begin{reglist}
\item [I2C\_RX\_LSB\_FIRST] 该位用于控制接收数据的存储模式。（读/写） \\1：从最低有效位开始接收数据；\\0：从最高有效位开始接收数据。
\item [I2C\_TX\_LSB\_FIRST] 该位用于控制发送数据的发送模式。（读/写）\\1：从最低有效位开始发送数据；\\0：从最高有效位开始发送数据。
\item [I2C\_TRANS\_START]置位该位，开始在 txfifo 中发送数据。（读/写）
\item [I2C\_MS\_MODE]置位该位，配置该模块为 I2C Master。清除该位，配置该模块为 I2C Slave。（读/写）
\item [I2C\_SAMPLE\_SCL\_LEVEL] 1：SCL 为低电平时，采样 SDA 数据；0：SCL 为高电平时，采样 SDA 数据。（读/写）
\item [I2C\_SCL\_FORCE\_OUT] 0：直接输出；  1：漏极开漏输出。（读/写）
\item [I2C\_SDA\_FORCE\_OUT] 0：直接输出；  1：漏极开漏输出。（读/写）
\end{reglist}\end{regdesc}
\end{register}


\begin{register}{H}{I2C\_SR\_REG}{0x0008}\label{regdesc:I2CSRREG}
\regfield{(reserved)}{1}{31}{0}%
\regfield{I2C\_SCL\_STATE\_LAST}{3}{28}{000}%
\regfield{(reserved)}{1}{27}{0}%
\regfield{I2C\_SCL\_MAIN\_STATE\_LAST}{3}{24}{000}%
\regfield{I2C\_TXFIFO\_CNT}{6}{18}{000000}%
\regfield{(reserved)}{4}{14}{0000}%
\regfield{I2C\_RXFIFO\_CNT}{6}{8}{000000}%
\regfield{(reserved)}{1}{7}{0}%
\regfield{I2C\_BYTE\_TRANS}{1}{6}{0}%
\regfield{I2C\_SLAVE\_ADDRESSED}{1}{5}{0}%
\regfield{I2C\_BUS\_BUSY}{1}{4}{0}%
\regfield{I2C\_ARB\_LOST}{1}{3}{0}%
\regfield{I2C\_TIME\_OUT}{1}{2}{0}%
\regfield{I2C\_SLAVE\_RW}{1}{1}{0}%
\regfield{I2C\_ACK\_REC}{1}{0}{0}%
\reglabel{Reset}\regnewline%
\begin{regdesc}\begin{reglist}
\item[I2C\_SCL\_STATE\_LAST] 存储产生 SCL 的状态机的值。（只读）\\
0：空闲；1：开始；2：下降沿；3：低电平；4：上升沿；5：高电平；6：停止
\item[I2C\_SCL\_MAIN\_STATE\_LAST] 存储 I2C 模块的状态机的值。（只读）\\
0：空闲；1：地址转换；2：ACK 地址；3：接收数据；4：发送数据；5：发送 ACK；6：等待 ACK
\item [I2C\_TXFIFO\_CNT]该字段存储了 RAM 中接收数据的字节数。（只读）
\item [I2C\_RXFIFO\_CNT]该字段表示需发送数据的字节数。（只读）
\item [I2C\_BYTE\_TRANS]当传输了一个字节时，该字段变为 1。（只读）
\item [I2C\_SLAVE\_ADDRESSED]当配置成 I2C Slave，主机发送的地址为从机地址时，该位翻转为高电平。（只读）
\item [I2C\_BUS\_BUSY]1：I2C 总线处于传输数据状态； 0：I2C 总线处于空闲状态。（只读）
\item [I2C\_ARB\_LOST]当 I2C 控制器不控制 SCL 线，该寄存器变为 1 。（只读）
\item [I2C\_TIME\_OUT]当 I2C 控制器接收一个 bit 的时间超时，该字段变为 1。（只读）
\item [I2C\_SLAVE\_RW]在从机模式中，1：主机在从机中读取数据；  0：主机向从机写入数据。 （只读）
\item [I2C\_ACK\_REC]该寄存器存储了接收到的 ACK 位的值。（只读）
\end{reglist}\end{regdesc}
\end{register}


\begin{register}{H}{I2C\_TO\_REG}{0x000c}\label{regdesc:I2CTOREG}
\regfield{(reserved)}{12}{20}{000000000000}%
\regfield{I2C\_TIME\_OUT\_REG}{20}{0}{00000000000000000000}%
\reglabel{Reset}\regnewline%
\begin{regdesc}\begin{reglist}
\item [I2C\_TIME\_OUT\_REG]当接收一个 bit 数据超过该值时，产生超时中断。（读/写）
\end{reglist}\end{regdesc}
\end{register}


\begin{register}{H}{I2C\_SLAVE\_ADDR\_REG}{0x0010}\label{regdesc:I2CSLAVEADDRREG}
\regfield{I2C\_SLAVE\_ADDR\_10BIT\_EN}{1}{31}{0}%
\regfield{(reserved)}{16}{15}{0000000000000000}%
\regfield{I2C\_SLAVE\_ADDR}{15}{0}{000000000000000}%
\reglabel{Reset}\regnewline%
\begin{regdesc}\begin{reglist}
\item [I2C\_SLAVE\_ADDR\_10BIT\_EN]I2C 为 Slave 时，该字段用于开启从机的 10-bit 寻址模式。（读/写）
\item [I2C\_SLAVE\_ADDR]当配置为 I2C Slave 时，该字段用于配置 Slave 地址。（读/写）
\end{reglist}\end{regdesc}
\end{register}


\begin{register}{H}{I2C\_RXFIFO\_ST\_REG}{0x0014}\label{regdesc:I2CRXFIFOSTREG}
\regfield{(reserved)}{12}{20}{000000000000}%
\regfield{I2C\_RXFIFO\_TXFIFO\_END\_ADDR}{5}{15}{00000}%
\regfield{I2C\_RXFIFO\_TXFIFO\_START\_ADDR}{5}{10}{00000}%
\regfield{I2C\_RXFIFO\_END\_ADDR}{5}{5}{00000}%
\regfield{I2C\_RXFIFO\_START\_ADDR}{5}{0}{00000}%
\reglabel{Reset}\regnewline%
\begin{regdesc}\begin{reglist}
\item [I2C\_TXFIFO\_END\_ADDR]如 nonfifo\_tx\_thres 所述，该位是最后发送数据的偏移地址。该值在每次 I2C\_TX\_SEND\_EMPTY\_INT 中断或 I2C\_TRANS\_COMPLETE\_INT 中断产生时都会更新。（只读）
\item [I2C\_TXFIFO\_START\_ADDR]如 nonfifo\_tx\_thres 所述，该位是最先发送数据的偏移地址。（只读）
\item [I2C\_RXFIFO\_END\_ADDR]如 nonfifo\_rx\_thres 所述，该位是最先接收数据的偏移地址。该值在每次 I2C\_RX\_REC\_FULL\_INT 中断或 I2C\_TRANS\_COMPLETE\_INT 中断产生时都会更新。（只读）
\item [I2C\_RXFIFO\_START\_ADDR]如 nonfifo\_rx\_thres 所述，该位是最后接收数据的偏移地址。（只读）
\end{reglist}\end{regdesc}
\end{register}


\begin{register}{H}{I2C\_FIFO\_CONF\_REG}{0x0018}\label{regdesc:I2CFIFOCONFREG}
\regfield{(reserved)}{6}{26}{000000}%
\regfield{I2C\_NONFIFO\_TX\_THRES}{6}{20}{{0x15}}%
\regfield{I2C\_NONFIFO\_RX\_THRES}{6}{14}{{0x15}}%
\regfield{(reserved)}{2}{12}{00}%
\regfield{I2C\_FIFO\_ADDR\_CFG\_EN}{1}{11}{0}%
\regfield{I2C\_NONFIFO\_EN}{1}{10}{0}%
\regfield{I2C\_TXFIFO\_EMPTY\_THRHD}{5}{5}{0}%
\regfield{I2C\_RXFIFO\_FULL\_THRHD}{5}{0}{0}%
\reglabel{Reset}\regnewline%
\begin{regdesc}\begin{reglist}
\item [I2C\_NONFIFO\_TX\_THRES]当 I2C 发送的数据多于 nonfifo\_tx\_thres 存储的字节，它将产生 tx\_send\_empty\_int\_raw 溢出中断，并更新发送数据的当前偏移地址。（读/写）
\item [I2C\_NONFIFO\_RX\_THRES]当 I2C 接收的数据多于 nonfifo\_tx\_thres 存储的字节，它将产生 rx\_send\_full\_int\_raw 溢出中断，并更新接收数据的当前偏移地址。（读/写） 
\item [I2C\_FIFO\_ADDR\_CFG\_EN]当该位为 1 时，可以按寄存器地址访问 Slave。（读/写）
\item [I2C\_NONFIFO\_EN]置位使能 APB nonfifo 访问方式。（读/写）
\item [I2C\_TXFIFO\_EMPTY\_THRHD]配置在直接访问模式下，TX FIFO 的水线。TX FIFO 计数值大于 I2C\_TXFIFO\_EMPTY\_THRHD [4:0] 时，I2C\_TXFIFO\_EMPTY\_INT\_RAW 位有效。（读/写）
\item [I2C\_RXFIFO\_FULL\_THRHD]配置在直接访问模式下，RX FIFO 的水线。RX FIFO 计数值大于 I2C\_RXFIFO\_FULL\_THRHD [4:0] 时，I2C\_RXFIFO\_FULL\_INT\_RAW 位有效。（读/写）
\end{reglist}\end{regdesc}
\end{register}


\begin{register}{H}{I2C\_INT\_RAW\_REG}{0x0020}\label{regdesc:I2CINTRAWREG}
\regfield{(reserved)}{19}{13}{0000000000000000000}%
\regfield{I2C\_TX\_SEND\_EMPTY\_INT\_RAW}{1}{12}{0}%
\regfield{I2C\_RX\_REC\_FULL\_INT\_RAW}{1}{11}{0}%
\regfield{I2C\_ACK\_ERR\_INT\_RAW}{1}{10}{0}%
\regfield{I2C\_TRANS\_START\_INT\_RAW}{1}{9}{0}%
\regfield{I2C\_TIME\_OUT\_INT\_RAW}{1}{8}{0}%
\regfield{I2C\_TRANS\_COMPLETE\_INT\_RAW}{1}{7}{0}%
\regfield{I2C\_MASTER\_TRAN\_COMP\_INT\_RAW}{1}{6}{0}%
\regfield{I2C\_ARBITRATION\_LOST\_INT\_RAW}{1}{5}{0}%
\regfield{I2C\_SLAVE\_TRAN\_COMP\_INT\_RAW}{1}{4}{0}%
\regfield{I2C\_END\_DETECT\_INT\_RAW}{1}{3}{0}%
\regfield{I2C\_RXFIFO\_OVF\_INT\_RAW}{1}{2}{0}%
\regfield{I2C\_TXFIFO\_EMPTY\_INT\_RAW}{1}{1}{0}%
\regfield{I2C\_RXFIFO\_FULL\_INT\_RAW}{1}{0}{0}%
\reglabel{Reset}\regnewline%
\begin{regdesc}\begin{reglist}
\item [I2C\_TX\_SEND\_EMPTY\_INT\_RAW] \hyperref[int:I2CTXSENDEMPTYINT]{I2C\_TX\_SEND\_EMPTY\_INT} 中断的原始中断状态位。（只读）
\item [I2C\_RX\_REC\_FULL\_INT\_RAW] \hyperref[int:I2CRXRECFULLINT]{I2C\_RX\_REC\_FULL\_INT} 中断的原始中断状态位。（只读）
\item [I2C\_ACK\_ERR\_INT\_RAW] \hyperref[int:I2CACKERRINT]{I2C\_ACK\_ERR\_INT} 中断的原始中断状态位。（只读）
\item [I2C\_TRANS\_START\_INT\_RAW] \hyperref[int:I2CTRANSSTARTINT]{I2C\_TRANS\_START\_INT} 中断的原始中断状态位。（只读）
\item [I2C\_TIME\_OUT\_INT\_RAW] \hyperref[int:I2CTIMEOUTINT]{I2C\_TIME\_OUT\_INT} 中断的原始中断状态位。（只读）
\item [I2C\_TRANS\_COMPLETE\_INT\_RAW] \hyperref[int:I2CTRANSCOMPLETEINT]{I2C\_TRANS\_COMPLETE\_INT} 中断的原始中断状态位。（只读）
\item [I2C\_MASTER\_TRAN\_COMP\_INT\_RAW] \hyperref[int:I2CMASTERTRANCOMPINT]{I2C\_MASTER\_TRAN\_COMP\_INT} 中断的原始中断状态位。（只读）
\item [I2C\_ARBITRATION\_LOST\_INT\_RAW] \hyperref[int:I2CARBITRATIONLOSTINT]{I2C\_ARBITRATION\_LOST\_INT} 中断的原始中断状态位。（只读）
\item [I2C\_SLAVE\_TRAN\_COMP\_INT\_RAW] \hyperref[int:I2CSLAVETRANCOMPINT]{I2C\_SLAVE\_TRAN\_COMP\_INT} 中断的原始中断状态位。（只读）
\item [I2C\_END\_DETECT\_INT\_RAW] \hyperref[int:I2CENDDETECTINT]{I2C\_END\_DETECT\_INT} 中断的原始中断状态位。（只读）
\item [I2C\_RXFIFO\_OVF\_INT\_RAW] \hyperref[int:I2CRXFIFOOVFINT]{I2C\_RXFIFO\_OVF\_INT} 中断的原始中断状态位。（只读）
\item [I2C\_TXFIFO\_EMPTY\_INT\_RAW] \hyperref[int:I2CTXFIFOEMPTYINT]{I2C\_TXFIFO\_EMPTY\_INT} 中断的原始中断状态位。（只读）
\item [I2C\_RXFIFO\_FULL\_INT\_RAW] \hyperref[int:I2CRXFIFOFULLINT]{I2C\_RXFIFO\_FULL\_INT} 中断的原始中断状态位。（只读）
\end{reglist}\end{regdesc}
\end{register}


\begin{register}{H}{I2C\_INT\_CLR\_REG}{0x0024}\label{regdesc:I2CINTCLRREG}
\regfield{(reserved)}{19}{13}{0000000000000000000}%
\regfield{I2C\_TX\_SEND\_EMPTY\_INT\_CLR}{1}{12}{0}%
\regfield{I2C\_RX\_REC\_FULL\_INT\_CLR}{1}{11}{0}%
\regfield{I2C\_ACK\_ERR\_INT\_CLR}{1}{10}{0}%
\regfield{I2C\_TRANS\_START\_INT\_CLR}{1}{9}{0}%
\regfield{I2C\_TIME\_OUT\_INT\_CLR}{1}{8}{0}%
\regfield{I2C\_TRANS\_COMPLETE\_INT\_CLR}{1}{7}{0}%
\regfield{I2C\_MASTER\_TRAN\_COMP\_INT\_CLR}{1}{6}{0}%
\regfield{I2C\_ARBITRATION\_LOST\_INT\_CLR}{1}{5}{0}%
\regfield{I2C\_SLAVE\_TRAN\_COMP\_INT\_CLR}{1}{4}{0}%
\regfield{I2C\_END\_DETECT\_INT\_CLR}{1}{3}{0}%
\regfield{I2C\_RXFIFO\_OVF\_INT\_CLR}{1}{2}{0}%
\regfield{I2C\_TXFIFO\_EMPTY\_INT\_CLR}{1}{1}{0}%
\regfield{I2C\_RXFIFO\_FULL\_INT\_CLR}{1}{0}{0}%
\reglabel{Reset}\regnewline%
\begin{regdesc}\begin{reglist}
\item [I2C\_TX\_SEND\_EMPTY\_INT\_CLR]置位该位用以清除 \hyperref[int:I2CTXSENDEMPTYINT]{I2C\_TX\_SEND\_EMPTY\_INT} 中断。（只写）
\item [I2C\_RX\_REC\_FULL\_INT\_CLR]置位该位用以清除 \hyperref[int:I2CRXRECFULLINT]{I2C\_RX\_REC\_FULL\_INT} 中断。（只写）
\item [I2C\_ACK\_ERR\_INT\_CLR]置位该位用以清除 \hyperref[int:I2CACKERRINT]{I2C\_ACK\_ERR\_INT} 中断。（只写）
\item [I2C\_TRANS\_START\_INT\_CLR]置位该位用以清除 \hyperref[int:I2CTRANSSTARTINT]{I2C\_TRANS\_START\_INT} 中断。（只写）
\item [I2C\_TIME\_OUT\_INT\_CLR]置位该位用以清除 \hyperref[int:I2CTIMEOUTINT]{I2C\_TIME\_OUT\_INT} 中断。（只写）
\item [I2C\_TRANS\_COMPLETE\_INT\_CLR]置位该位用以清除 \hyperref[int:I2CTRANSCOMPLETEINT]{I2C\_TRANS\_COMPLETE\_INT} 中断。（只写）
\item [I2C\_MASTER\_TRAN\_COMP\_INT\_CLR]置位该位用以清除 \hyperref[int:I2CMASTERTRANCOMPINT]{I2C\_MASTER\_TRAN\_COMP\_INT} 中断。（只写）
\item [I2C\_ARBITRATION\_LOST\_INT\_CLR]置位该位用以清除 \hyperref[int:I2CARBITRATIONLOSTINT]{I2C\_ARBITRATION\_LOST\_INT} 中断。（只写）
\item [I2C\_SLAVE\_TRAN\_COMP\_INT\_CLR]置位该位用以清除 \hyperref[int:I2CSLAVETRANCOMPINT]{I2C\_SLAVE\_TRAN\_COMP\_INT} 中断。（只写）
\item [I2C\_END\_DETECT\_INT\_CLR]置位该位用以清除 \hyperref[int:I2CENDDETECTINT]{I2C\_END\_DETECT\_INT} 中断。（只写）
\item [I2C\_RXFIFO\_OVF\_INT\_CLR]置位该位用以清除 \hyperref[int:I2CRXFIFOOVFINT]{I2C\_RXFIFO\_OVF\_INT} 中断。（只写）
\item [I2C\_TXFIFO\_EMPTY\_INT\_CLR]置位该位用以清除 \hyperref[int:I2CTXFIFOEMPTYINT]{I2C\_TXFIFO\_EMPTY\_INT} 中断。（只写）
\item [I2C\_RXFIFO\_FULL\_INT\_CLR]置位该位用以清除 \hyperref[int:I2CRXFIFOFULLINT]{I2C\_RXFIFO\_FULL\_INT} 中断。（只写）
\end{reglist}\end{regdesc}
\end{register}


\begin{register}{H}{I2C\_INT\_ENA\_REG}{0x0028}\label{regdesc:I2CINTENAREG}
\regfield{(reserved)}{19}{13}{0000000000000000000}%
\regfield{I2C\_TX\_SEND\_EMPTY\_INT\_ENA}{1}{12}{0}%
\regfield{I2C\_RX\_REC\_FULL\_INT\_ENA}{1}{11}{0}%
\regfield{I2C\_ACK\_ERR\_INT\_ENA}{1}{10}{0}%
\regfield{I2C\_TRANS\_START\_INT\_ENA}{1}{9}{0}%
\regfield{I2C\_TIME\_OUT\_INT\_ENA}{1}{8}{0}%
\regfield{I2C\_TRANS\_COMPLETE\_INT\_ENA}{1}{7}{0}%
\regfield{I2C\_MASTER\_TRAN\_COMP\_INT\_ENA}{1}{6}{0}%
\regfield{I2C\_ARBITRATION\_LOST\_INT\_ENA}{1}{5}{0}%
\regfield{I2C\_SLAVE\_TRAN\_COMP\_INT\_ENA}{1}{4}{0}%
\regfield{I2C\_END\_DETECT\_INT\_ENA}{1}{3}{0}%
\regfield{I2C\_RXFIFO\_OVF\_INT\_ENA}{1}{2}{0}%
\regfield{I2C\_TXFIFO\_EMPTY\_INT\_ENA}{1}{1}{0}%
\regfield{I2C\_RXFIFO\_FULL\_INT\_ENA}{1}{0}{0}%
\reglabel{Reset}\regnewline%
\begin{regdesc}\begin{reglist}
\item [I2C\_TX\_SEND\_EMPTY\_INT\_ENA]中断使能位，用于 \hyperref[int:I2CTXSENDEMPTYINT]{I2C\_TX\_SEND\_EMPTY\_INT} 中断。（读/写）
\item [I2C\_RX\_REC\_FULL\_INT\_ENA]中断使能位，用于 \hyperref[int:I2CRXRECFULLINT]{I2C\_RX\_REC\_FULL\_INT} 中断。（读/写）
\item [I2C\_ACK\_ERR\_INT\_ENA]中断使能位，用于 \hyperref[int:I2CACKERRINT]{I2C\_ACK\_ERR\_INT} 中断。（读/写）
\item [I2C\_TRANS\_START\_INT\_ENA]中断使能位，用于 \hyperref[int:I2CTRANSSTARTINT]{I2C\_TRANS\_START\_INT} 中断。（读/写）
\item [I2C\_TIME\_OUT\_INT\_ENA]中断使能位，用于 \hyperref[int:I2CTIMEOUTINT]{I2C\_TIME\_OUT\_INT} 中断。（读/写）
\item [I2C\_TRANS\_COMPLETE\_INT\_ENA]中断使能位，用于 \hyperref[int:I2CTRANSCOMPLETEINT]{I2C\_TRANS\_COMPLETE\_INT} 中断。（读/写）
\item [I2C\_MASTER\_TRAN\_COMP\_INT\_ENA]中断使能位，用于 \hyperref[int:I2CMASTERTRANCOMPINT]{I2C\_MASTER\_TRAN\_COMP\_INT} 中断。（读/写）
\item [I2C\_ARBITRATION\_LOST\_INT\_ENA]中断使能位，用于 \hyperref[int:I2CARBITRATIONLOSTINT]{I2C\_ARBITRATION\_LOST\_INT} 中断。（读/写）
\item [I2C\_SLAVE\_TRAN\_COMP\_INT\_ENA]中断使能位，用于 \hyperref[int:I2CSLAVETRANCOMPINT]{I2C\_SLAVE\_TRAN\_COMP\_INT} 中断。（读/写）
\item [I2C\_END\_DETECT\_INT\_ENA]中断使能位，用于 \hyperref[int:I2CENDDETECTINT]{I2C\_END\_DETECT\_INT} 中断。（读/写）
\item [I2C\_RXFIFO\_OVF\_INT\_ENA]中断使能位，用于 \hyperref[int:I2CRXFIFOOVFINT]{I2C\_RXFIFO\_OVF\_INT} 中断。（读/写）
\item [I2C\_TXFIFO\_EMPTY\_INT\_ENA]中断使能位，用于 \hyperref[int:I2CTXFIFOEMPTYINT]{I2C\_TXFIFO\_EMPTY\_INT} 中断。（读/写）
\item [I2C\_RXFIFO\_FULL\_INT\_ENA]中断使能位，用于 \hyperref[int:I2CRXFIFOFULLINT]{I2C\_RXFIFO\_FULL\_INT} 中断。（读/写）
\end{reglist}\end{regdesc}
\end{register}


\begin{register}{H}{I2C\_INT\_STATUS\_REG}{0x002c}\label{regdesc:I2CINTSTATUSREG}
\regfield{(reserved)}{19}{13}{0000000000000000000}%
\regfield{I2C\_TX\_SEND\_EMPTY\_INT\_ST}{1}{12}{0}%
\regfield{I2C\_RX\_REC\_FULL\_INT\_ST}{1}{11}{0}%
\regfield{I2C\_ACK\_ERR\_INT\_ST}{1}{10}{0}%
\regfield{I2C\_TRANS\_START\_INT\_ST}{1}{9}{0}%
\regfield{I2C\_TIME\_OUT\_INT\_ST}{1}{8}{0}%
\regfield{I2C\_TRANS\_COMPLETE\_INT\_ST}{1}{7}{0}%
\regfield{I2C\_MASTER\_TRAN\_COMP\_INT\_ST}{1}{6}{0}%
\regfield{I2C\_ARBITRATION\_LOST\_INT\_ST}{1}{5}{0}%
\regfield{I2C\_SLAVE\_TRAN\_COMP\_INT\_ST}{1}{4}{0}%
\regfield{I2C\_END\_DETECT\_INT\_ST}{1}{3}{0}%
\regfield{I2C\_RXFIFO\_OVF\_INT\_ST}{1}{2}{0}%
\regfield{I2C\_TXFIFO\_EMPTY\_INT\_ST}{1}{1}{0}%
\regfield{I2C\_RXFIFO\_FULL\_INT\_ST}{1}{0}{0}%
\reglabel{Reset}\regnewline%
\begin{regdesc}\begin{reglist}
\item [I2C\_TX\_SEND\_EMPTY\_INT\_ST]隐蔽中断状态位，用于 \hyperref[int:I2CTXSENDEMPTYINT]{I2C\_TX\_SEND\_EMPTY\_INT} 中断。（只读） 
\item [I2C\_RX\_REC\_FULL\_INT\_ST]隐蔽中断状态位，用于 \hyperref[int:I2CRXRECFULLINT]{I2C\_RX\_REC\_FULL\_INT}中断。（只读）
\item [I2C\_ACK\_ERR\_INT\_ST]隐蔽中断状态位，用于 \hyperref[int:I2CACKERRINT]{I2C\_ACK\_ERR\_INT} 中断。（只读）
\item [I2C\_TRANS\_START\_INT\_ST]隐蔽中断状态位，用于 \hyperref[int:I2CTRANSSTARTINT]{I2C\_TRANS\_START\_INT} 中断。（只读）
\item [I2C\_TIME\_OUT\_INT\_ST]隐蔽中断状态位，用于 \hyperref[int:I2CTIMEOUTINT]{I2C\_TIME\_OUT\_INT} 中断。（只读）
\item [I2C\_TRANS\_COMPLETE\_INT\_ST]隐蔽中断状态位，用于 \hyperref[int:I2CTRANSCOMPLETEINT]{I2C\_TRANS\_COMPLETE\_INT} 中断。（只读）
\item [I2C\_MASTER\_TRAN\_COMP\_INT\_ST]隐蔽中断状态位，用于 \hyperref[int:I2CMASTERTRANCOMPINT]{I2C\_MASTER\_TRAN\_COMP\_INT} 中断。（只读）
\item [I2C\_ARBITRATION\_LOST\_INT\_ST]隐蔽中断状态位，用于 \hyperref[int:I2CARBITRATIONLOSTINT]{I2C\_ARBITRATION\_LOST\_INT} 中断。（只读）
\item [I2C\_SLAVE\_TRAN\_COMP\_INT\_ST]隐蔽中断状态位，用于 \hyperref[int:I2CSLAVETRANCOMPINT]{I2C\_SLAVE\_TRAN\_COMP\_INT} 中断。（只读）
\item [I2C\_END\_DETECT\_INT\_ST]隐蔽中断状态位，用于\hyperref[int:I2CENDDETECTINT]{I2C\_END\_DETECT\_INT} 中断。（只读）
\item [I2C\_RXFIFO\_OVF\_INT\_ST]隐蔽中断状态位，用于 \hyperref[int:I2CRXFIFOOVFINT]{I2C\_RXFIFO\_OVF\_INT} 中断。（只读）
\item [I2C\_TXFIFO\_EMPTY\_INT\_ST]隐蔽中断状态位，用于 \hyperref[int:I2CTXFIFOEMPTYINT]{I2C\_TXFIFO\_EMPTY\_INT} 中断。（只读）
\item [I2C\_RXFIFO\_FULL\_INT\_ST]隐蔽中断状态位，用于 \hyperref[int:I2CRXFIFOFULLINT]{I2C\_RXFIFO\_FULL\_INT} 中断。（只读）
\end{reglist}\end{regdesc}
\end{register}


\begin{register}{H}{I2C\_SDA\_HOLD\_REG}{0x0030}\label{regdesc:I2CSDAHOLDREG}
\regfield{(reserved)}{22}{10}{0000000000000000000000}%
\regfield{I2C\_SDA\_HOLD\_TIME}{10}{0}{0000000000}%
\reglabel{Reset}\regnewline%
\begin{regdesc}\begin{reglist}
\item [I2C\_SDA\_HOLD\_TIME]
用于配置 SCL 下降沿后 SDA 保持不变的 APB 时钟周期数。（读/写）
\end{reglist}\end{regdesc}
\end{register}


\begin{register}{H}{I2C\_SDA\_SAMPLE\_REG}{0x0034}\label{regdesc:I2CSDASAMPLEREG}
\regfield{(reserved)}{22}{10}{0000000000000000000000}%
\regfield{I2C\_SDA\_SAMPLE\_TIME}{10}{0}{0000000000}%
\reglabel{Reset}\regnewline%
\begin{regdesc}\begin{reglist}
\item [I2C\_SDA\_SAMPLE\_TIME] 用于配置 I2C 采样 SDA 的时间，基于 APB 时钟周期数。（读/写）
\end{reglist}\end{regdesc}
\end{register}


\begin{register}{H}{I2C\_SCL\_HIGH\_PERIOD\_REG}{0x0038}\label{regdesc:I2CSCLHIGHPERIODREG}
\regfield{(reserved)}{18}{14}{000000000000000000}%
\regfield{I2C\_SCL\_HIGH\_PERIOD}{14}{0}{00000000000000}%
\reglabel{Reset}\regnewline%
\begin{regdesc}\begin{reglist}
\item [I2C\_SCL\_HIGH\_PERIOD]I2C 为 Master 时，用于配置使 SCL 维持在高电平的 APB 时钟周期数。（读/写）
\end{reglist}\end{regdesc}
\end{register}


\begin{register}{H}{I2C\_SCL\_START\_HOLD\_REG}{0x0040}\label{regdesc:I2CSCLSTARTHOLDREG}
\regfield{(reserved)}{22}{10}{0000000000000000000000}%
\regfield{I2C\_SCL\_START\_HOLD\_TIME}{10}{0}{0000001000}%
\reglabel{Reset}\regnewline%
\begin{regdesc}\begin{reglist}
\item [I2C\_SCL\_START\_HOLD\_TIME] 用于配置 SDA 下降沿和 SCL 下降沿之间的时间，用于 START 状态，基于 APB 时钟周期数。（读/写）
\end{reglist}\end{regdesc}
\end{register}


\begin{register}{H}{I2C\_SCL\_RSTART\_SETUP\_REG}{0x0044}\label{regdesc:I2CSCLRSTARTSETUPREG}
\regfield{(reserved)}{22}{10}{0000000000000000000000}%
\regfield{I2C\_SCL\_RSTART\_SETUP\_TIME}{10}{0}{0000001000}%
\reglabel{Reset}\regnewline%
\begin{regdesc}\begin{reglist}
\item [I2C\_SCL\_RSTART\_SETUP\_TIME] 用于配置 SCL 的上升沿和 SDA 的下降沿之间的时间，用于 RESTART 状态，基于 APB 时钟周期数。（读/写）
\end{reglist}\end{regdesc}
\end{register}


\begin{register}{H}{I2C\_SCL\_STOP\_HOLD\_REG}{0x0048}\label{regdesc:I2CSCLSTOPHOLDREG}
\regfield{(reserved)}{18}{14}{000000000000000000}%
\regfield{I2C\_SCL\_STOP\_HOLD\_TIME}{14}{0}{00000000000000}%
\reglabel{Reset}\regnewline%
\begin{regdesc}\begin{reglist}
\item [I2C\_SCL\_STOP\_HOLD\_TIME] 用于配置 STOP 状态后的延迟，基于 APB 时钟周期数。（读/写）
\end{reglist}\end{regdesc}
\end{register}


\begin{register}{H}{I2C\_SCL\_STOP\_SETUP\_REG}{0x004C}\label{regdesc:I2CSCLSTOPSETUPREG}
\regfield{(reserved)}{22}{10}{0000000000000000000000}%
\regfield{I2C\_SCL\_STOP\_SETUP\_TIME}{10}{0}{0000000000}%
\reglabel{Reset}\regnewline%
\begin{regdesc}\begin{reglist}
\item [I2C\_SCL\_STOP\_SETUP\_TIME] 用于配置 SCL 上升沿和 SDA 上升沿之间的时间，基于 APB 时钟周期数。（读/写）
\end{reglist}\end{regdesc}
\end{register}


\begin{register}{H}{I2C\_SCL\_FILTER\_CFG\_REG}{0x0050}\label{regdesc:I2CSCLFILTERCFGREG}
\regfield{(reserved)}{28}{4}{0000000000000000000000000000}%
\regfield{I2C\_SCL\_FILTER\_EN}{1}{3}{1}%
\regfield{I2C\_SCL\_FILTER\_THRES}{3}{0}{000}%
\reglabel{Reset}\regnewline%
\begin{regdesc}\begin{reglist}
\item [I2C\_SCL\_FILTER\_EN]该位为 SCL 的滤波使能位。（读/写）
\item [I2C\_SCL\_FILTER\_THRES]基于 APB 时钟周期数，当 SCL 输入信号上的脉冲宽度小于该寄存器的值时， I2C 控制器 将过滤该脉冲。（读/写） 
\end{reglist}\end{regdesc}
\end{register}


\begin{register}{H}{I2C\_SDA\_FILTER\_CFG\_REG}{0x0054}\label{regdesc:I2CSDAFILTERCFGREG}
\regfield{(reserved)}{28}{4}{0000000000000000000000000000}%
\regfield{I2C\_SDA\_FILTER\_EN}{1}{3}{1}%
\regfield{I2C\_SDA\_FILTER\_THRES}{3}{0}{000}%
\reglabel{Reset}\regnewline%
\begin{regdesc}\begin{reglist}
\item [I2C\_SDA\_FILTER\_EN]该位为 SDA 的滤波使能位。（读/写）
\item [I2C\_SDA\_FILTER\_THRES]基于 APB 时钟周期数，当 SDA 输入信号上的脉冲宽度小于此寄存器的值时， I2C 控制器 将过滤该脉冲。（读/写） 
\end{reglist}\end{regdesc}
\end{register}


\begin{register}{H}{I2C\_COMD\textcolor{red}{\it{n}}\_REG (\textcolor{red}{\it{n}}: 0-15)}{0x58+4*\textcolor{red}{\it{n}}}\label{regdesc:I2CCOMDnREG}
\regfield{I2C\_COMMAND\textcolor{red}{\it{n}}\_DONE}{1}{31}{0}%
\regfield{(reserved)}{17}{14}{00000000000000000}%
\regfield{I2C\_COMMAND\textcolor{red}{\it{n}}}{14}{0}{00000000000000}%
\reglabel{Reset}\regnewline%
\begin{regdesc}\begin{reglist}
\item [I2C\_COMMAND\textcolor{red}{\it{n}}\_DONE]在 I2C 主机模式下完成 command \textcolor{red}{\it{n}} 时，该位翻转为高电平。（读/写）
\item [I2C\_COMMAND\textcolor{red}{\it{n}}] command \textcolor{red}{\textcolor{red}{\it{n}}}  的内容。它包括三个部分：（读/写）\\
op\_code 是命令，0：START；1： WRITE；2： READ；3：STOP；4：END。\\
Byte\_num 代表了需要被发送或接收的字节数。\\
ack\_check\_en， ack\_exp 和 ack 用于控制 ACK 位。详情请参考 \hyperref[para:i2ccmd]{I2C cmd 结构}。
\end{reglist}\end{regdesc}
\end{register}



\end{document}
